\section{Cote 68: exponentials and rearrangements}\label{sec:part-c}

\subsection{The corrected exponential correspondence}
Throughout this subsection, $K$ is a commutative $\Q$-algebra and $A$ is a
commutative unital $K$-algebra. Nilpotent for an endomorphism means nilpotent
with one bound valid on the whole module; locally nilpotent means that a
bound may depend on the vector. These meanings must be kept separate in
cote 68, scans 039 and 044--046. We impose no Noetherian or finite generation
hypothesis.

\begin{theorem}[Proved; polynomial exponential and logarithm]\label{C-exp-log}
The maps
\[
 D\longmapsto \exp D=\sum_{j\geq0}\frac{D^j}{j!},\qquad
 u\longmapsto \log u=\sum_{j\geq1}\frac{(-1)^{j+1}}j(u-1)^j
\]
are mutually inverse bijections between nilpotent $K$-derivations of $A$
and $K$-algebra automorphisms $u$ for which $u-1$ is nilpotent in
$\operatorname{End}_K(A)$. Both sums are finite. The nilpotency indices of
$D$ and $\exp D-1$ coincide.
\end{theorem}
\begin{proof}
For a derivation, induction gives
$D^n(ab)=\sum_{i=0}^n\binom niD^i(a)D^{n-i}(b)$.
Consequently the finite exponential preserves multiplication and the unit;
$\exp(-D)$ is its inverse. Polynomial identities in $\Q[T]/(T^m)$ show
$\log(\exp T)=T$ and $\exp(\log(1+T))=1+T$ when $T^m=0$.
These identities prove that the displayed operations are inverse as
operations on nilpotent linear endomorphisms. It remains to prove that
$\log u$ is a derivation.

Put $N=u-1$ and choose $m$ with $N^m=0$. Define a polynomial family of
$K$-linear maps $P_t:A\to A[t]$ by
\[
 P_t=\sum_{j=0}^{m-1}\binom tjN^j.
\]
For every nonnegative integer $q$, the binomial theorem gives $P_q=u^q$;
hence $P_q(ab)=P_q(a)P_q(b)$ and $P_q(1)=1$. A polynomial with coefficients
in a $\Q$-vector space vanishing at all nonnegative integers is zero:
for a polynomial of degree at most $d$, its values at $0,\ldots,d$ determine
its coefficients through an invertible rational Vandermonde matrix.
Therefore $P_t(ab)=P_t(a)P_t(b)$ in $A[t]$. Differentiating at $t=0$ and
using
\[
 \left.\frac{d}{dt}\binom tj\right|_{t=0}
 =\frac{(-1)^{j-1}}j\quad(j\geq1)
\]
gives the Leibniz identity for $\log u$. It kills $K$ because $N$ kills $K$.
Finally $\exp D-1=Dq(D)$ for a polynomial $q$ with $q(0)=1$;
$q(D)$ is invertible and commutes with $D$. Thus the two nilpotency
indices coincide. The same argument applies to $\log(1+N)=Nr(N)$.
\end{proof}

\begin{theorem}[Proved; C.1(a), the struck condition (iii$'$)]\label{C-filtered}
Let $I\subset A$ be an ideal with $I^m=0$, where $m\geq1$. Exponential is a
bijection from
\[
 \mathfrak d_I=
 \{D\in\operatorname{Der}_K(A):D(I^k)\subset I^{k+1}
                         \text{ for every }k\geq0\}
\]
to
\[
 U_I=\{u\in\operatorname{Aut}_K(A):u(I)=I,
                 \ \operatorname{gr}_I(u)=1\}.
\]
Its inverse is the polynomial logarithm. Here $I^0=A$ and
$\operatorname{gr}_I(A)=\bigoplus_{k\geq0}I^k/I^{k+1}$.
This is the valid replacement of the last corollary of 68/046.
\end{theorem}
\begin{proof}
If $D\in\mathfrak d_I$, then $D^j(I^k)\subset I^{k+j}$ and $D^m=0$.
Every term of $\exp D-1$ raises the filtration at least once, so
$\exp D\in U_I$. Conversely, if $u\in U_I$, then
$N=u-1$ satisfies $N(I^k)\subset I^{k+1}$, hence $N^m=0$.
Theorem~\ref{C-exp-log} proves that $\log u$ is a $K$-derivation, and
its polynomial formula proves that it raises the filtration. The inverse
identities were proved in the same theorem.
\end{proof}

We use the differential-operator convention of
\cite[Definition 10.133.1 and Lemma 10.133.2, tags 09CI and 09CJ]{StacksDiff}.
It is the affine commutator convention of EGA IV, 16.8.8.
The following calculation supplies the exact specialization needed here.

\begin{theorem}[Proved; C.1(b), differential order and conormal criterion]\label{C-differential}
For $u\in\operatorname{Aut}_K(A)$ let
$J=(u(a)-a:a\in A)\subset A$. For every $n\geq0$,
\[
 u\text{ has differential order at most }n
 \quad\Longleftrightarrow\quad J^{n+1}=0.
\]
Suppose more generally that $I^m=0$ and $u$ induces the identity on $A/I$.
Then $u(I)=I$ and it induces an $(A/I)$-linear automorphism
$\gamma:I/I^2\to I/I^2$. The following are equivalent:
\begin{enumerate}
\item $u=\exp D$ for a nilpotent $K$-derivation $D$ of $A$;
\item $u-1$ is nilpotent on the $K$-module $A$;
\item $\gamma-1$ is nilpotent on the $(A/I)$-module $I/I^2$.
\end{enumerate}
If $(\gamma-1)^r=0$, then one valid bound is
\[
 (u-1)^{m+(r-1)m(m-1)/2}=0.
\]
In particular this characterizes exactly which differential automorphisms
are exponentials, by taking $I=J$. The other differential automorphisms
are exactly those with non-unipotent action on $J/J^2$, where unipotent
here requires a uniform nilpotence bound.
\end{theorem}
\begin{proof}
Write $M_a$ for multiplication by $a$ on $A$. Direct calculation gives
\[
 [u,M_a]=M_{u(a)-a}u,
 \qquad
 [\cdots[[u,M_{a_0}],M_{a_1}],\ldots,M_{a_n}]
 =M_{\prod_{j=0}^n(u(a_j)-a_j)}u.
\]
The latter operator is zero exactly when the product is zero, as follows
by evaluating at $1$. The recursive definition of differential order
therefore gives the first equivalence.

The equation $\pi u=\pi$, for $\pi:A\to A/I$, implies $u(I)\subset I$.
If $u(x)\in I$, then $\pi(x)=\pi u(x)=0$, proving equality $u(I)=I$.
For $a\in A$ and $x\in I$, the difference
$u(a)u(x)-a u(x)$ lies in $I^2$; thus $\gamma$ is $(A/I)$-linear.
The equivalence of (1) and (2) is Theorem~\ref{C-exp-log}, and (2)
implies (3) by passage to an invariant subquotient.

For the converse put $R=A/I$ and $V=I/I^2$. Multiplication induces a
surjective $R$-linear map $V^{\otimes_R k}\to I^k/I^{k+1}$ that intertwines
$\gamma^{\otimes k}$ with the action of $u$. On the tensor product, put
$N_j=1\otimes\cdots\otimes(\gamma-1)\otimes\cdots\otimes1$.
The $N_j$ commute and $N_j^r=0$. Therefore the ideal they generate in their
commutative endomorphism algebra has power $k(r-1)+1$ equal to zero:
every monomial of that degree has some exponent at least $r$.
Since $\gamma^{\otimes k}-1=\prod_j(1+N_j)-1$ lies in this ideal,
\[
 (u-1)^{k(r-1)+1}(I^k)\subset I^{k+1}\quad(k\geq1).
\]
Also $(u-1)(A)\subset I$. Composing these inclusions for
$k=0,\ldots,m-1$, with exponent $1$ for $k=0$, proves the asserted bound.
The proof used no flatness, finite generation, or Noetherian assumption.
\end{proof}

\begin{proposition}[Refuted; 68/045--046, and exact replacement]\label{C-other-autos}
A differential automorphism need not be an exponential of a nilpotent
derivation. On $A=\Q[\epsilon]/(\epsilon^2)$, the automorphism
$u(\epsilon)=2\epsilon$ has differential order one and is not such an
exponential. The assertion in 68/045 that $u$ automatically acts identically
on $\operatorname{gr}_J(A)$ is also false.

For any $\Q$-vector space $V$, the square-zero algebra $\Q\oplus V$ has
all $\Q$-algebra automorphisms of the form $1\oplus T$, with
$T\in\operatorname{GL}_{\Q}(V)$. All are differential of order at most one;
precisely those for which $T-1$ is uniformly nilpotent are exponentials of
nilpotent derivations. Thus no finite dimensional or semisimple description
of the remaining automorphisms is available without additional hypotheses.
\end{proposition}
\begin{proof}
In the first example $J=(\epsilon)$, so $J^2=0$ and
Theorem~\ref{C-differential} applies. On $J/J^2$ the operator is
multiplication by $2$, hence is not unipotent. This also disproves the
asserted associated graded identity.

The multiplication in $\Q\oplus V$ is
$(q,v)(q',v')=(qq',qv'+q'v)$. Its nilradical is $V$, so every automorphism
preserves $V$, fixes $\Q$, and is exactly $1\oplus T$ for an invertible
linear map $T$. Conversely each such map preserves this multiplication.
Its difference ideal is contained in $V$ and has square zero.
The criterion follows from Theorem~\ref{C-exp-log} directly on
$\Q\oplus V$. In particular take a direct sum of finite Jordan blocks
with unbounded sizes for $T-1$. This gives a locally unipotent differential
automorphism which is not uniformly unipotent and is not the exponential
of any uniformly nilpotent derivation. Finally on
$\Q[\epsilon]/(\epsilon^3)$ the derivation
$D=\epsilon\,d/d\epsilon$ maps $A$ into $(\epsilon)$, but
$D^j(\epsilon)=\epsilon$ for every $j\geq1$. This checks the failure of
68/046; the hypothesis repaired in Theorem~\ref{C-filtered} excludes it.
\end{proof}

\begin{theorem}[Proved; C.1(c), locally nilpotent intertwining]\label{C-intertwine}
Let $M,N$ be arbitrary $K$-modules, let $u\in\operatorname{End}_K(M)$ and
$v\in\operatorname{End}_K(N)$ be locally nilpotent, and let $f:M\to N$ be
$K$-linear. Then
\[
 f\exp(u)=\exp(v)f\quad\Longleftrightarrow\quad fu=vf.
\]
The exponentials are pointwise finite, mutually invertible linear maps.
\end{theorem}
\begin{proof}
For $m\in M$, the cyclic $K$-submodule generated by
$m,u(m),u^2(m),\ldots$ is killed by a power of $u$.
All formal exponential and logarithm identities therefore hold on this
submodule as finite polynomial identities. In particular
$U=\exp u$ satisfies that $U-1$ is locally nilpotent and
$\log U=u$ pointwise. The same holds for $V=\exp v$.
If $fU=Vf$, induction gives $f(U-1)^j=(V-1)^jf$. For a fixed $m$, choose
one bound large enough to truncate both logarithm series at $m$ and at
$f(m)$. Summing the preceding identities with coefficients
$(-1)^{j+1}/j$ proves $fu(m)=vf(m)$. Conversely $fu=vf$ implies the
corresponding identity for every power and then for the pointwise finite
exponential sums. This argument does not assert that the induced operator
$h\mapsto vh-hu$ on the entire module $\operatorname{Hom}_K(M,N)$ is
locally nilpotent; that assertion is unnecessary and can fail.
\end{proof}

\subsection{What changes for arbitrary module functors}
A $K$-module functor $\mathcal M$ assigns an $R$-module $\mathcal M(R)$ to
every commutative $K$-algebra $R$, with semilinear transition maps. Its
endomorphisms are natural transformations that are $R$-linear on every
$\mathcal M(R)$. For a $K$-algebra $R$, let
$\mathscr E(R)=\operatorname{End}_{R}(\mathcal M_R)$ be the ring of such
endomorphisms after restricting to $R$-algebras. This is itself a module
functor. Put $\mathcal O(R)=R$ and $\mathcal N(R)=\operatorname{Nil}(R)$.
A group action in the next statements is a natural group homomorphism
into the units of $\mathscr E$; this specifies the target meant by
$\operatorname{Aut}_K(\mathcal M)$ in 68/044.

\begin{theorem}[Proved; C.1(d), exact logarithmic criteria]\label{C-functor-criterion}
For any $K$-module functor $\mathcal M$:
\begin{enumerate}
\item The ordinary exponential map
\[
 \{u\in\mathscr E(K):u\text{ nilpotent}\}
 \longrightarrow\operatorname{Hom}_{\mathrm{gr}}(\mathcal O,\mathscr E^\times)
\]
is injective. Its image consists exactly of actions $\rho$ satisfying both
of the following explicit conditions: there is a single integer $m$ with
$(\rho_R(\lambda)-1)^m=0$ for every $R$ and $\lambda\in R$; and their
polynomial logarithms define an $\mathcal O$-linear natural map
$L:\mathcal O\to\mathscr E$. Thus this exponential map is bijective
exactly when every action satisfies these two conditions.
\item The formal exponential map
\[
 \mathscr E(K)\longrightarrow
 \operatorname{Hom}_{\mathrm{gr}}(\mathcal N,\mathscr E^\times)
\]
has image exactly the formal actions $\rho$ for which each
$\rho_R(\lambda)-1$ is nilpotent in $\mathscr E(R)$ and the natural map
$L_R(\lambda)=\log\rho_R(\lambda)$ extends to an $\mathcal O$-linear
natural map $\mathcal O\to\mathscr E$. Nilpotence bounds in this assertion
may depend on $(R,\lambda)$. Its fibres, when nonempty, are translates of
\[
 H(\mathcal M)=\{w\in\mathscr E(K):
 \lambda w_R=0\text{ for every }R\text{ and every }
 \lambda\in\operatorname{Nil}(R)\}.
\]
Consequently formal exponential is bijective exactly when
$H(\mathcal M)=0$ and every formal action meets the displayed logarithm
and extension conditions.
\end{enumerate}
These criteria apply to every module functor; they require no
quasicoherence assumption.
\end{theorem}
\begin{proof}
For nilpotent $u$, evaluation of its ordinary exponential at
$1\in K$ gives $\exp u$. The inverse polynomial logarithm recovers $u$,
proving injectivity. Its values have a uniform nilpotence bound because
$\exp(\lambda u_R)-1$ is $u_R$ times an invertible polynomial, and
its logarithm is $\lambda u_R$, an $\mathcal O$-linear map.
Conversely an $\mathcal O$-linear natural map
$L:\mathcal O\to\mathscr E$ is determined by $u=L_K(1)$: naturality and
$R$-linearity give $L_R(\lambda)=\lambda u_R$. The assumed uniform
nilpotence makes $u$ nilpotent; exponentiating the logarithm gives
$\rho_R(\lambda)=\exp(\lambda u_R)$.

For the formal case, $\lambda^m=0$ implies $(\lambda u_R)^m=0$ for any
$u$. Hence the exponential has the claimed logarithm, which extends by
$\lambda\mapsto\lambda u_R$. In the converse direction an extension of
$L$ is again determined by $u=L_K^{\mathrm{ext}}(1)$, and finite exp--log
identities recover the action. No nilpotence condition on $u$ is required.
If $u,v$ have the same formal exponential, taking finite logarithms at
each nilpotent scalar gives $\lambda(u_R-v_R)=0$. Conversely this equality
makes their exponentials equal. This proves the fibre statement and both
bijectivity criteria. The additive group law ensures that the logarithm
is additive whenever defined: the two commuting nilpotent operators
$\rho_R(\lambda)-1$ and $\rho_R(\mu)-1$ generate a nilpotent commutative
subalgebra, in which the ordinary polynomial logarithm identity applies.
The extension condition imposes the additional scalar linearity needed
for an exponential; additivity alone does not impose it.
\end{proof}

\begin{proposition}[Refuted; universal functor bijectivity in 68/044]\label{C-functor-counter}
The formal exponential need not be injective, and neither exponential
map need be surjective, for arbitrary module functors.
\end{proposition}
\begin{proof}
For $K=\Q$, take $\mathcal M(R)=R_{\mathrm{red}}$ with its quotient
$R$-module structure. Nilpotent scalars kill this functor, including
after every base extension. Thus $H(\mathcal M)=\mathscr E(K)$, which
contains the nonzero identity. Every formal exponential is the identity
action, so formal exponential is not injective.

For failures of surjectivity, take
$\mathcal M(R)=R\oplus\Omega_{R/\Q}$ and define
\[
 \rho_R(\lambda)(a,\omega)=(a,\omega+a\,d\lambda).
\]
These maps are $R$-linear, are natural in $R$, and satisfy
$\rho(\lambda+\mu)=\rho(\lambda)\rho(\mu)$; their inverses are
$\rho(-\lambda)$. Moreover $(\rho(\lambda)-1)^2=0$.
The ordinary action has $\rho(1)=1$. If it were $e_u$ for a nilpotent
natural endomorphism $u$, polynomial logarithm at $1$ would give $u=0$,
whereas $\rho_{\Q[t]}(t)(1,0)=(1,dt)\ne(1,0)$.

The restricted formal action is not an exponential either. Let
$R=\Q[\epsilon]/(\epsilon^2)$. Then
$\Omega_{R/\Q}=\Q\,d\epsilon$ and
$\epsilon\Omega_{R/\Q}=0$. For any natural endomorphism $u$,
$\exp(\epsilon u_R)=1+\epsilon u_R$, whose difference from the identity
has zero $\Omega$-component. But
$(\rho_R(\epsilon)-1)(1,0)=(0,d\epsilon)\ne0$.
All algebraic hypotheses and naturality assertions are explicit in these
constructions.
\end{proof}

\begin{theorem}[Known, with complete derivation; quasicoherent case]\label{C-qcoh}
Let $M$ be any $K$-module and $\mathcal M(R)=R\otimes_K M$. Formal actions
of $\widehat{\mathbb G}_a=\mathcal N$ on $\mathcal M$ are uniquely the
$\bar e_u$ for $u\in\operatorname{End}_K(M)$. Ordinary actions of
$\mathbb G_a=\mathcal O$ are uniquely the $e_u$ for locally nilpotent
endomorphisms $u$ of $M$. Consequently the ordinary map with domain only
uniformly nilpotent endomorphisms is bijective exactly when every locally
nilpotent endomorphism of $M$ is uniformly nilpotent. Finite generation of
$M$ suffices for this last property but is not asserted to be necessary.
The ordinary classification is stated in
\cite[Chapter VIII, Example 2.1 and footnote 1, p.~113]{MilneAGS}.
\end{theorem}
\begin{proof}
For a formal action, evaluate at the universal nilpotent element $t$ in
$K[t]/(t^{n+1})$. Applied to $1\otimes m$, the resulting element has a
unique expansion
$\sum_{j=0}^{n}t^j\otimes u_j(m)$. Compatibility under reduction in $n$
gives $K$-linear maps $u_j:M\to M$ independent of $n$. Naturality under
$t\mapsto\lambda$ and $R$-linearity determine the action from these maps
on every base algebra. The identity and group laws, evaluated in
$K[s,t]/(s^{n+1},t^{n+1})$, give
\[
 u_0=1,\qquad \binom{i+j}{i}u_{i+j}=u_i u_j.
\]
Putting $i=1$ and inducting yields $u_j=u_1^j/j!$, which proves existence
and uniqueness of the formal description.

For an ordinary action, use the same construction in $K[t]$.
For each $m$, the image is a polynomial in $t$, so $u_j(m)=0$ for all
large $j$. The same identities give $u_j(m)=u_1^j(m)/j!$; hence $u_1$ is
locally nilpotent. Conversely a locally nilpotent endomorphism remains
locally nilpotent after every base change: an element of $R\otimes M$
is a finite sum of pure tensors and one takes the maximum of finitely many
nilpotence bounds. Its exponential therefore defines a natural action.
Uniqueness also follows from the coefficient of $t$. If $M$ has finitely
many generators, the maximum of their individual bounds kills $M$,
proving the final sufficient condition. For $M=K[x]$ over $K=\Q$, the
endomorphism $d/dx$ is locally nilpotent but has no uniform bound, giving
the correction required in 68/039. Finite generation is not necessary:
take $K=\Q[x]$ and $M=\bigoplus_{n\geq0}K/(x-n)$. Distinct summands
have no nonzero homomorphisms between them, because $x-n$ acts invertibly
on $K/(x-j)$ for $j\ne n$. Thus every endomorphism is diagonal, with
one rational scalar on each summand. Local nilpotence forces every scalar
to be zero. Yet $M$ is not finitely generated, because finitely many
generators have support in only finitely many summands. This refines the
phrase ``only for finitely generated $M$'' in the supplied study.
\end{proof}

\subsection{Rearrangement: definitions and the missing implication}
Let $G$ be any group, without a commutativity or countability assumption.
We write $\R^{(G)}_+$ for the nonnegative real functions with finite support,
$\delta_a$ for the point mass at $a$, and
\[
 (f*g)(x)=\sum_{y\in G}f(y)g(y^{-1}x),\qquad
 \check f(x)=f(x^{-1}).
\]
Two such functions are equimeasurable, written $f\sim g$, when their
positive values, with multiplicity, agree. Write $f^\downarrow_j$ for the
decreasing list, padded with zeros, and define
\[
 |f,g|=\sum_{j\geq1}f^\downarrow_jg^\downarrow_j,
 \qquad f'\prec f\ \Longleftrightarrow\
 \sum_{j=1}^k(f')^\downarrow_j\leq\sum_{j=1}^kf^\downarrow_j
 \text{ for every }k.
\]
When $G$ is finite, it suffices to take $1\leq k\leq|G|$.
A pair is \emph{pr\'etass\'e} if $\|f*g\|_\infty=|f,g|$; it is
\emph{tass\'e} if $f'*g'\prec f*g$ whenever $f'\prec f$ and $g'\prec g$.
These are the mathematical notions on 68/020--022. The original scan 017
has been inspected independently; the image-reading report supplies the
names, and no result below depends on the uncertain spelling.

\begin{lemma}[Proved; rearrangement and equality]\label{C-rearrangement}
For $f,g\in\R^{(G)}_+$,
\[
 \sup_{f'\sim f}\sum_x f'(x)g(x)=|f,g|,
 \qquad \|f*g\|_\infty\leq|f,g|.
\]
An inner product achieves its rearrangement maximum precisely when every
positive superlevel set of one factor is comparable by inclusion with
every positive superlevel set of the other. Consequently $(f,g)$ is
pr\'etass\'e precisely when there exists $s\in G$ such that the superlevel
sets of $f$ and $\delta_s*\check g$ together form a chain.
Moreover $f'\prec f$ is equivalent to $|f',h|\leq|f,h|$ for every
$h\in\R^{(G)}_+$.
\end{lemma}
\begin{proof}
Both functions have finitely many positive values. Expand each as a
positive linear combination of the indicator functions of its distinct
positive superlevel sets. Every intersection of sets of sizes $a,b$ has
size at most $\min(a,b)$, with equality exactly when they are comparable
by inclusion. Arranging both lists of values in the same decreasing
order attains all these upper bounds at once. The positive coefficients
in the expansions prove the stated equality criterion as well. This is
the finite rearrangement inequality, also recorded in
\cite[Chapter X, \S10.2, Theorem 368]{HLP}.

For convolution,
$(f*g)(s)=\langle f,\delta_s*\check g\rangle$, and the second factor is
equimeasurable with $g$. The convolution has finite support, so its
maximum is attained. This proves both convolution assertions.
Finally, with $h^\downarrow_{r+1}=0$, summation by parts gives
\[
 |f,h|=\sum_{k=1}^{r}
   (h^\downarrow_k-h^\downarrow_{k+1})
   \sum_{j=1}^kf^\downarrow_j.
\]
This proves one direction of the last equivalence. The reverse direction
uses an indicator function of each permitted cardinality $k$ as $h$.
\end{proof}

For clarity we state the typed version of the draft axioms on 68/017.
There is a group $\widetilde G$, a homomorphism
$\rho:\widetilde G\to G$, and for each $a\in\widetilde G$ a family
$\mathcal A_a$ of finite subsets of $G$. The cone $C_a$ is the set of finite
nonnegative linear combinations of their indicators. The axioms used are:
\begin{enumerate}
\item[(a)] $\mathcal A_a$ is a chain under inclusion;
\item[(b)] $A=\rho(a)A^{-1}$ for $A\in\mathcal A_a$;
\item[(c)] $1_A*1_B\in C_{ab}$ for $A\in\mathcal A_a$, $B\in\mathcal A_b$;
\item[(d)] for $a,b$ there is $s\in G$ such that
      $\mathcal A_a\cup s\mathcal A_b$ is a chain.
\end{enumerate}
The page also assumes $i:{}_2G\simeq\ker\rho$ and writes, in a separate
clause labelled (d),
$\mathcal A_{i(t)a}=t\mathcal A_a$ for $t\in{}_2G$.
The supplied transcription has $B\in\mathcal B(\widetilde h)$ in (c),
although only $\mathcal A$ is defined. We state the type-correct reading
$\mathcal B=\mathcal A$. The counterexample below also applies if a second
family is allowed, by taking it equal to the first. There is no completed
fullness or nonemptiness requirement on this page; the final bracket is
unfinished. We make this distinction explicit in the classification below.

\begin{theorem}[Proved; what the chain axioms imply]\label{C-chain-true}
Under (a)--(d), $S=\bigcup_{a\in\widetilde G}C_a$ is convolution stable,
and every pair in $S\times S$ is pr\'etass\'e. This conclusion does not
require the additional kernel and torsion equivariance assumptions.
\end{theorem}
\begin{proof}
Bilinearity and (c) show $C_a*C_b\subset C_{ab}$.
The positive superlevel sets of a finite nonnegative combination of
indicators in a chain are themselves members of that chain, apart from
the empty set. Axiom (b) also gives $f=\delta_{\rho(a)}*\check f$ for
$f\in C_a$. For $f\in C_a$ and $g\in C_b$, choose $s$ using (d).
The superlevel sets of $f$ and $\delta_s*g$ form a chain. As
$g=\delta_{\rho(b)}*\check g$, Lemma~\ref{C-rearrangement} applies with
translation $s\rho(b)$ and proves the assertion.
\end{proof}

\begin{theorem}[Refuted; C.2(a) and Lead L8]\label{C-K-counter}
The axioms on 68/017 do not imply Condition $K$ on 68/021, even with all
positive cardinalities represented, nonempty families, and the exact
specified two-torsion kernel. This failure already occurs on $G=\Z/4\Z$.
\end{theorem}
\begin{proof}
Use additive notation, set $\widetilde G=\Z/8\Z$, and let
$\rho(k)=k\bmod4$ and $i(t)=2t\bmod8$ for $t\in{}_2G=\{0,2\}$.
Put $A=\{3,0,1\}$ and $B=\{0,1\}$ in $G$. For $a\bmod4$ define
\[
 \mathcal A_{2a}=\{\{a\},A+a,G\},\qquad
 \mathcal A_{2a+1}=\{B+a,G\}.
\]
Each family is a chain and each of its sets is invariant under
$x\mapsto k-x$ in fibre $k$, proving (a) and (b). The two-torsion
identification is exact, and adding $i(2)=4$ to the index translates the
family by $2$, proving the extra equivariance condition.

For (c), translation reduces all products to the following complete
list, with subsets identified with their indicators:
\[
 \begin{array}{c|ccc}
 *&\delta_0&A&B\\ \hline
 \delta_0&\delta_0&A&B\\
 A&A&2G+\delta_0&G+B\\
 B&B&G+B&(A+1)+\delta_1
 \end{array}
 \qquad G*1_D=|D|G.
\]
Every right-hand side belongs to its required fibre cone. Translation of
both inputs translates the output by the sum, so this verifies every
atom product. To check (d), translate each even fibre to
$\{\{0\},A,G\}$ and each odd fibre to $\{B,G\}$.
These two chains interlace because $\{0\}\subset B\subset A\subset G$.
Two fibres of the same parity can be made equal by a translation.
The families also represent all cardinalities $1,2,3,4$.

Now let $f=1_B$ and $h=1_{\{0,2\}}$. Both have decreasing value list
$(1,1,0,0)$, so $h\sim f$ and in particular $h\prec f$. Listed at
$0,1,2,3$, the convolution powers are
\[
 \begin{array}{c|cc}
       &\text{square}&\text{cube}\\ \hline
 f &(1,2,1,0)&(1,3,3,1)\\
 h &(2,0,2,0)&(4,0,4,0).
 \end{array}
\]
Thus the three-factor instance of Condition $K$ would assert $4\leq3$.
Already the two-factor weak-majorization claim fails at the sum of the
two largest values, which is $4$ for $h*h$ and $3$ for $f*f$.
The supplied exact program checks all $400$ atom products, all $64$
alignment pairs, all $16$ torsion-equivariance pairs, and the displayed
convolutions. The hand table supplies the proof independently of that
finite certificate.
\end{proof}

\begin{theorem}[Proved; exact replacement for Condition $K$]\label{C-K-replacement}
Let $S\subset\R^{(G)}_+$ be convolution stable and suppose that for every
positive integer $k\leq|G|$ it contains an indicator $1_{A_k}$ with
$|A_k|=k$. The following are equivalent:
\begin{enumerate}
\item Every pair in $S\times S$ is tass\'e.
\item For every $n\geq2$, $f_i\in S$ and $f_i'\prec f_i$ imply
      $f_1'*\cdots*f_n'\prec f_1*\cdots*f_n$.
\item Condition $K$ holds: under those hypotheses,
      $\|f_1'*\cdots*f_n'\|_\infty
          \leq\|f_1*\cdots*f_n\|_\infty$.
\end{enumerate}
For the implication from (3) to (1), it suffices to assume only the
three-factor instances with third factor $1_{A_k}$. Thus the missing
hypothesis in Lead L8 is a genuine convolution-majorization property;
chain closure and pairwise optimal suprema alone do not supply it.
\end{theorem}
\begin{proof}
Assume (1). Induct on $n$. If the first $n-1$ primed factors have
convolution $h'\prec h$, then $h\in S$ by convolution stability, and
applying the tass\'e property to $(h,f_n)$ proves the next step.
The implication (2)$\Rightarrow$(3) is the first partial-sum inequality.

Assume the indicated three-factor cases of (3). Fix $f,g\in S$,
$f'\prec f$, $g'\prec g$, and $k$. Choose an indicator $1_{A_k'}$
equimeasurable with $1_{A_k}$ whose support, after inversion, consists of
$k$ points realizing the $k$ largest values of $f'*g'$, padded by zeros if
necessary. Then
\[
 \begin{split}
 \sum_{j=1}^k(f'*g')^\downarrow_j
 &=(f'*g'*1_{A_k'})(e)\\
 &\leq\|f*g*1_{A_k}\|_\infty
 \leq |f*g,1_{A_k}|
 =\sum_{j=1}^k(f*g)^\downarrow_j.
 \end{split}
\]
All comparisons are valid also for infinite $G$, since every support is
finite. This proves (1).
\end{proof}

\paragraph{The precise classical source.}
Gabriel's original paper \cite[\S1, pp.~32--34]{Gabriel1932C} already
distinguishes two orders for finite nonnegative sequences on $\Z$:
$0,1,-1,2,-2,\ldots$ and its reflection. Denote the corresponding
rearrangements by $a^+$ and $a^-$. A symmetric multiset in his sense
has a symmetric decreasing representative $c^*$. His Theorem~1,
p.~34, proves
\[
 (a*b*c)(0)\leq(a^+*b^-*c^*)(0)
\]
when the values of $c$ form such a multiset; his Theorem~2, p.~38,
allows any additional symmetric factors. The opposite orders of the
first two arbitrary factors are essential: pp.~33--34 explicitly
reject the proposed inequality using the same one-sided order in
every factor. This is a theorem for sequences on $\Z$ with its stated
symmetry hypotheses. It does not supply Condition~$K$ for the
arbitrary group cone in 68/017. The counterexample and exact
majorization replacement above are independent of that classical
result.

\subsection{Existence of systems and the two-torsion indexing}

\begin{theorem}[Proved under the literal interpretation; C.2(b)]\label{C-literal-class}
Interpret ${}_2G$ as the subset $T=\{t\in G:t^2=e\}$ with the multiplication
inherited from $G$, and retain exactly the draft axioms just stated,
without adding cardinality completeness. For a finite or arbitrary
discrete group $G$, data satisfying these axioms and
$i:T\simeq\ker\rho$ exist if and only if $T$ is a subgroup of $G$,
equivalently if and only if every two involutions commute. The construction
can have a surjective $\rho$ and at least one nonzero cone. In particular,
within the abelian category every finite and every discrete group admits
a system in this literal sense.
\end{theorem}
\begin{proof}
For the stipulated isomorphism of groups with the inherited multiplication
to make sense, $T$ must be a subgroup. If $s,t\in T$, then
$(st)^2=e$ is equivalent to $st=ts$. This proves the equivalence of the
two group-theoretic conditions.

Conversely, suppose $T$ is a subgroup; it is then elementary abelian.
Set $\widetilde G=T\times G$, let $\rho(t,g)=g$, and put
$i(t)=(t,e)$. Define
\[
 \mathcal A_{(t,g)}=
 \begin{cases}
 \{\{t\}\},&g=e,\\
 \varnothing,&g\ne e.
 \end{cases}
\]
These are families of finite sets, and the fibre $(e,e)$ has a nonzero
cone. Chains and the symmetry condition are immediate: in a nonempty
fibre the centre is $e$ and $t=t^{-1}$.
The only atom products to check are
$\delta_t*\delta_s=\delta_{ts}$ in the kernel fibres, exactly as required.
Multiplying the index by $i(u)$ replaces $\{t\}$ by $\{ut\}$, giving
kernel equivariance. Two nonempty fibres become the same singleton family
after translation by $ts^{-1}$; alignment with an empty family imposes no
condition. Thus all axioms hold and $\rho$ is surjective.
The proof is uniform for finite and infinite groups.
\end{proof}

This answers the existence question for the axioms actually written; it
is not a classification of cardinality-complete, faithful, or maximal
systems. None of those additional adjectives appears as a completed
hypothesis on 68/017. If one intends such an additional requirement, it
must be named as a different problem. In the nonabelian category where
$T$ is not a subgroup, the notation $i:{}_2G\simeq\ker\rho$ is not a
well-defined group isomorphism with inherited multiplication, rather than
an obstruction supplied by an unproved classification theorem.

\begin{proposition}[Proved; the extension is extra data, not forced]\label{C-extension}
The chain, symmetry, convolution, and alignment axioms do not determine
the kernel of $\rho$. Even after the additional requirement
$\ker\rho\simeq{}_2G$ is imposed, they do not determine the extension up
to group isomorphism, including when all possible positive set cardinalities occur.
\end{proposition}
\begin{proof}
Without the kernel requirement, take any homomorphism $\rho$ and let
every family consist only of the empty set; all four axioms hold for
arbitrary kernels. This proves logical independence of the requirement.

For nonuniqueness with the requirement, take $G=\Z/2\Z$. One system has
$\widetilde G=\Z/4\Z$, $\rho(k)=k\bmod2$, and $i(t)=2t$.
Its families are
\[
 \mathcal A_0=\{\{0\},G\},\quad
 \mathcal A_2=\{\{1\},G\},\quad
 \mathcal A_1=\mathcal A_3=\{G\}.
\]
A second system has $\widetilde G=G\times G$, $\rho(s,t)=s$,
$i(t)=(0,t)$, and families
\[
 \mathcal A_{(0,t)}=\{\{t\},G\},\qquad
 \mathcal A_{(1,t)}=\{G\}.
\]
In both cases singletons convolve by their group product, $G*\delta_t=G$,
and $G*G=2G$, so every required cone condition holds. Symmetry, alignment,
and kernel equivariance follow from the displayed families. Cardinalities
one and two are represented. The index groups have different exponents,
four and two, so they are not isomorphic. These finite checks are also
included in the attached computation.
\end{proof}

\subsection{The symmetry question}

\begin{theorem}[Proved; C.2(c), exact group criterion]\label{C-symmetry}
For a group $G$, the pr\'etass\'e relation on all of
$\R^{(G)}_+\times\R^{(G)}_+$ is symmetric if and only if $G$ is abelian.
If both $f$ and $g$ are individually equal to their inversions up to a
left translation, their particular pr\'etass\'e relation is symmetric on
any group, as asserted on 68/021.
\end{theorem}
\begin{proof}
If $G$ is abelian, convolution commutes, proving the first implication.
Suppose $a,b\in G$ do not commute. Set
\[
 f=\delta_e+2\delta_a,\qquad
 g=\delta_b+2\delta_{a^{-1}b}.
\]
Both functions have decreasing positive value list $(2,1)$, so $|f,g|=5$.
Multiplication gives
\[
 \begin{split}
 f*g&=5\delta_b+2\delta_{a^{-1}b}+2\delta_{ab},\\
 g*f&=\delta_b+2\delta_{ba}+2\delta_{a^{-1}b}
                         +4\delta_{a^{-1}ba}.
 \end{split}
\]
In the first line, neither off-centre point equals $b$; if those points
coincide their combined coefficient is only $4$. Thus $\|f*g\|_\infty=5$.
In the second line $a^{-1}ba\ne b$, by noncommutativity. It cannot equal
$ba$ or $a^{-1}b$, since either equality would force $a=e$.
The remaining two coefficients of size $2$ can at most combine to $4$,
and neither combines with the point $b$. Hence $\|g*f\|_\infty=4$.
This proves failure of symmetry for every nonabelian group, finite or
infinite, with no subgroup classification assumption.

For the final assertion, write $f=\delta_r*\check f$ and
$g=\delta_s*\check g$. Lemma~\ref{C-rearrangement} says that
$(f,g)$ is pr\'etass\'e if and only if the superlevel sets of $f$ and some
left translate of $g$ form a chain. Translating this whole chain back
reverses the roles of $f,g$. Applying the same lemma and the assumed
symmetry of $f$ proves the equivalence with $(g,f)$.
\end{proof}

\begin{proposition}[Refuted by a finite certificate; symmetry of tass\'e pairs]\label{C-packed-asymmetry}
On $S_3$ the tass\'e relation also need not be symmetric. With
$a=(01)$ and $b=(12)$ in Theorem~\ref{C-symmetry}, the pair $(f,g)$ is
tass\'e and $(g,f)$ is not.
\end{proposition}
\begin{proof}
The decreasing convolution lists are $(5,4,0,0,0,0)$ and
$(4,2,2,1,0,0)$. Every $f'\prec f$ and $g'\prec g$ has
$\|f'*g'\|_\infty\leq|f',g'|\leq|f,g|=5$ by
Lemma~\ref{C-rearrangement}, used twice with summation by parts.
Also $\|f'\|_1\leq3$ and $\|g'\|_1\leq3$, so the total convolution mass
is at most $9$. These two bounds prove
$f'*g'\prec(5,4,0,0,0,0)$ for every permitted pair.
The reversed pair is not even pr\'etass\'e and hence cannot be tass\'e:
choosing equimeasurable rearrangements realizing inner product $5$
violates its first partial sum. The integer certificate independently
enumerates all $30\times30=900$ equimeasurable rearrangement pairs.
\end{proof}

\paragraph{Scope of the requested determinant source.}
The finite-free axioms on cote 68/056 are the degree-zero part of the
determinant formalism in Knudsen--Mumford
\cite[Chapter I, pp.~20--25, Definition~1 and Theorem~1]{KnudsenMumford1976C}.
Their full system additionally treats complexes, quasi-isomorphisms and
base change; we do not identify that larger system with the three
conditions on the sheet. The Koszul interchange sign is
$(-1)^{rs}$ for ranks $r,s$. The printed exponent $r+s$ on their p.~20
is a typographical error: interchanging two rank-one summands changes
their exterior product by $-1$. The detailed source comparison and
splitting-independence calculation are recorded in the source-audit note.
